% !TEX root = ../main.tex
\section{Results}
\label{sec:results}

This chapter reports the generalisation of every model to the held-out 30-scene
suite, organised by research question. Section~\ref{sec:rq_baseline} establishes
that the task is solvable. Section~\ref{sec:rq1} addresses whether asymmetric
self-play succeeds under the single-GPU budget (RQ1). Section~\ref{sec:rq2}
isolates the setup dimension that separates success from failure (RQ2).
Section~\ref{sec:rq3} identifies the underlying mechanism (RQ3).
Section~\ref{sec:cross_env_results} reports the second-simulator result.

Three success measures are reported and kept distinct throughout, and all are
defined in Section~\ref{sec:metrics}. The
\emph{held-out scene success rate} is computed on the fixed 30-scene suite that no
model saw during training, and is the measure of generalisation this thesis reports
as its headline. The \emph{within-distribution scene success rate} is computed on
goals drawn from the model's own training distribution, and its gap against the
held-out rate is the mechanism of Section~\ref{sec:rq3}. The \emph{training-time
success rate} is an auxiliary per-push rate computed on the goals the models
trained on. Unless stated otherwise, headline models report 5 seeds at 528
environments; ablations and scale sweeps report 3 seeds; confidence intervals are
95\,\% across seeds. The held-out numbers in this chapter follow the evaluation
protocol defined in Section~\ref{sec:val_protocol}, and every comparison reported
here is already decisive under that protocol.
Table~\ref{tab:main_results} collects the principal numbers.

\begin{table}[H]
    \centering
    \small
    \begin{tabular}{l c c c c}
        \hline
        \textbf{Model} & \textbf{Scene SR} & \textbf{Trial SR} & \textbf{Pos-only} & \textbf{Seeds} \\
        \hline
        PBRS (potential-based)   & \textbf{79.3\,\% $\pm$ 3.5\,\%} & 70.2\,\% & 100\,\% & 5 \\
        Curriculum               & \textbf{80.0\,\% $\pm$ 4.1\,\%} & 70.2\,\% & 100\,\% & 5 \\
        ASP-disc (position-only) & 7.3\,\% $\pm$ 1.8\,\%  & 5.8\,\% & --- & 5 \\
        TASP-disc                & 6.7\,\% $\pm$ 0.0\,\%  & 4.9\,\% & --- & 5 \\
        ASP-dPose (outcome)      & 0.0\,\%                & 0.0\,\% & 0\,\%  & 5 \\
        TASP-dPose (time-based)  & 0.0\,\%                & 0.0\,\% & 0\,\%  & 4 \\
        Bob-penalty              & 0.0\,\%                & 0.0\,\% & 0\,\%  & 4 \\
        \hline
    \end{tabular}
    \caption{Held-out generalisation of the headline models at 528 environments.
    Scene SR counts a scene solved if any of its attempts passes the combined gate;
    trial SR is the per-attempt rate. Pos-only is the scene SR on the ten
    position-only scenes. Values are seed means with 95\,\% confidence intervals
    across seeds, under the evaluation protocol of Section~\ref{sec:val_protocol}.}
    \label{tab:main_results}
\end{table}

\subsection{Baseline}
\label{sec:rq_baseline}

The single-agent models establish that the task is solvable, which is the
precondition for interpreting any self-play result as a property of the method
rather than the task. PBRS reaches a held-out scene success rate of
79.3\,\% ($\pm$ 3.5 percentage points, 5 seeds) at 528 environments, and the
curriculum model reaches 80.0\,\% ($\pm$ 4.1\,pp, 5 seeds). On the ten
position-only scenes both reach 100\,\%; on the twenty scenes that also require
an orientation change they reach 69\,\% and 70\,\%. The single-agent models
therefore fail only on rotation-requiring scenes, and never because of position.
Their trial success rate of 70.2\,\% shows that the policy is reliable as well as
capable: a scene is solved in the majority of attempts, not in a single lucky one.
The per-scene errors of Figure~\ref{fig:per_scene_errors} show where the
remaining failures concentrate: position error stays below tolerance on every
scene, and the unsolved scenes are those whose orientation error exceeds
the 0.2\,rad threshold.

The single-agent models show no generalisation gap. Their training goals are
sampled uniformly over the workspace, so the within-distribution evaluation of
Section~\ref{sec:metrics} draws goals from the same distribution as the held-out
suite, and the two success rates coincide. The 79.3\,\% held-out result is
therefore also the single-agent's within-distribution competence. This
distinguishes the single-agent models from the self-play models of
Section~\ref{sec:rq3}, whose within-distribution success collapses on the
held-out suite.

\begin{figure}[H]
    \centering
    \includegraphics[width=\textwidth,height=0.5\textheight,keepaspectratio]{per_scene_errors_3models.png}
    \caption{Per-scene generalisation errors, seed means across 5 seeds
    ($\pm$ 1 standard deviation). Top row: position error; bottom row:
    orientation error. Columns: PBRS, Curriculum, and ASP-dPose. Green bars mark
    scenes solved in the majority of seeds; the dashed line marks the combined-gate
    threshold. The single-agent models leave only rotation-dominant scenes
    unsolved, while ASP-dPose solves none.}
    \label{fig:per_scene_errors}
\end{figure}

\subsection{Research Question 1: ASP performance under the single-GPU budget}
\label{sec:rq1}

The self-play models do not generalise to the held-out suite at all. ASP-dPose
and TASP-dPose each solve zero of the 30 scenes, in every seed (5 and 4 seeds
respectively). The Bob-penalty model is likewise at zero. This is not a
marginal underperformance against the 79\,\% of the single-agent model; it is a
complete failure to transfer.

The two Alice reward formulations are indistinguishable at this level. The
outcome-based reward of \cite{OpenAI2021AsymmetricSF} and the time-based reward
of \cite{Sukhbaatar2018IntrinsicMotivation} both produce zero held-out successes,
so the reward-formulation dimension cannot be resolved on the T-block: both variants
fail entirely. The same holds for the symmetric Bob-penalty variant. The failure
is one of transfer, not of learning. On goals the goal-proposing agent itself
produces, the self-play policy solves 61.4\,\% of disc scenes and 11.7\,\% of
T-block scenes (5 seeds each, Section~\ref{sec:rq3}), and none of that
competence reaches the held-out suite. The
comparisons that can be made at this level are between the within-distribution
behaviour of Section~\ref{sec:rq3} and the held-out result, where a consistent
picture emerges.

The disc models, whose goal constrains position alone, are marginally above
zero. ASP-disc solves 7.3\,\% of its 30 disc scenes ($\pm$ 1.8\,pp, 5 seeds) and
TASP-disc 6.7\,\% ($\pm$ 0.0\,pp, 5 seeds). The couple of scenes each disc model
solves are the easiest of the suite. This places the disc contrast of
Section~\ref{sec:rq2_coupling} in its proper light: self-play fails on both
objects, and fails completely on the coupled-goal T-block.

\begin{figure}[H]
    \centering
    \includegraphics[width=\textwidth,height=0.45\textheight,keepaspectratio]{val_multi_combined.png}
    \caption{Held-out scene success rates for the headline models at 528
    environments. The single-agent models solve roughly four scenes in five; the
    disc self-play models solve fewer than one in ten; the T-block self-play
    models solve none.}
    \label{fig:val_multi}
\end{figure}

\subsection{Research Question 2: Setup dimensions}
\label{sec:rq2}

This section reports the single-dimension isolations that decompose the self-play
result. Each isolates one factor that differs between the successful prior
settings and the present setup.

\subsubsection{Goal Coupling}
\label{sec:rq2_coupling}

The disc self-play models isolate the coupled goal as a single variable. On the
rotationally symmetric disc, whose goal constrains position alone, self-play
solves 7.3\,\% of held-out scenes. On the T-block, whose goal couples position
and orientation, the same mechanism solves 0.0\,\%. The pose, the push primitive,
the Bob reward, and the self-play loop are identical between the two variants;
only the object and the success criterion differ.

The within-distribution evaluation of Section~\ref{sec:rq3} measures the coupling
directly. On goals the goal-proposing agent itself generates, the disc model
solves 61.4\,\% of scenes ($\pm$ 2.1\,pp, 5 seeds) while the T-block model
solves 11.7\,\% ($\pm$ 1.5\,pp, 5 seeds). The coupled gate therefore reduces
within-distribution competence by roughly 50 percentage points before any
transfer is required. The held-out suite then collapses both rates, the disc to
7.3\,\% and the T-block to zero. The failure thus separates into two effects of
comparable size. The coupled gate accounts for the within-distribution drop from
61.4\,\% to 11.7\,\%; the generalisation collapse accounts for the drop from each
within-distribution rate to its held-out rate, roughly 54 percentage points on
the disc. Neither effect alone explains the full result, and the second is the
subject of Section~\ref{sec:rq3}.

The single-agent models show the same gate effect in a solvable setting. On the
ten position-only scenes they reach 100\,\%; on the twenty scenes that also
demand an orientation change they reach 69--70\,\%. The coupled gate is hard even
for the model that solves the task, which is why it compounds self-play's
generalisation failure decisively.

\subsubsection{Reward Formulation}
\label{sec:rq2_reward}

The two Alice reward formulations produce identical held-out outcomes: zero
scenes solved for both ASP-dPose and TASP-dPose. The self-regulating property of
the time-based reward of \cite{Sukhbaatar2018IntrinsicMotivation} does not confer
a measurable advantage in this setup, and neither reward transfers. The
reward-formulation dimension is therefore not the factor that
separates the successful prior settings from this one: both formulations fail
equally here.

\subsubsection{Compute Scale}
\label{sec:rq2_scale}

The two self-play subjects were trained at 256, 528, and 2{,}048 parallel
environments with matched total experience, so that the per-update batch size
varies while the total pushes stay fixed. On the held-out suite ASP-dPose solves
0.0\,\% of scenes at 256 and 528 environments (3 and 5 seeds), and 1.1\,\%
($\pm$ 4.8\,pp, 3 seeds) at 2{,}048, where a single seed solves a single scene.
TASP-dPose solves 0.0\,\% at all three counts. An eightfold increase in the
parallel batch does not lift the self-play result above the noise floor. Compute
scale is therefore ruled out as the cause of the underperformance, in the
strongest sense available: the curve is not merely flat, it is flat at zero.

\subsubsection{Imitation Signal}
\label{sec:rq2_abc}

The imitation term introduced by \cite{OpenAI2021AsymmetricSF}, Alice
Behavioural Cloning at $\beta = 0.5$, is instrumented through the fraction of
demonstrated actions whose probability ratio falls outside the clipping interval.
For every ABC-enabled model this clip fraction is \textbf{0.000} at every
training iteration: the goal-proposing agent's exploratory pushes lie so far from
the goal-solving agent's policy that the clipping zeroes the gradient, and the
imitation signal never flows.

The ablation shows the signal's absence has no measurable effect because there
is nothing left to preserve. Removing ABC changes the held-out scene SR of ASP-dPose from
0.0\,\% to 0.0\,\% (3 seeds), and the same holds for TASP-dPose and for the disc
models (6.7\,\% with and without). The positive control on the disc, where the
task is easiest and the imitation term would be most likely to help, also shows
no effect. The practical consequence is that the imitation term contributes
nothing in this discrete macro-action domain; the clip-saturation measurement
explains why.

\subsubsection{Goal Representation}
\label{sec:rq2_encoder}

Removing the GoalEncoder from Bob and providing the full object-and-goal
observation directly to the policy encoder changes the held-out scene SR of
ASP-dPose from 0.0\,\% to 0.0\,\% (3 seeds). The eight-dimensional latent
compression is not the factor that limits self-play, consistent with the disc
isolation, where the same encoder is present and self-play retains its marginal
generalisation.

\subsection{Research Question 3: Mechanism}
\label{sec:rq3}

The mechanism of the self-play failure is the gap between the within-distribution
competence and the held-out performance.
Table~\ref{tab:rq3} and Figure~\ref{fig:train_heldout} place the two measures
side by side.

\begin{table}[H]
    \centering
    \small
    \begin{tabular}{l c c c}
        \hline
        \textbf{Model} & \textbf{Within-dist SR} & \textbf{Held-out SR} & \textbf{Seeds} \\
        \hline
        ASP-disc     & 61.4\,\% $\pm$ 2.1\,\%  & 7.3\,\%          & 5 \\
        ASP-dPose    & 11.7\,\% $\pm$ 1.5\,\%  & 0.0\,\%          & 5 \\
        \hline
    \end{tabular}
    \caption{Within-distribution against held-out scene success for the two
    self-play subjects. Within-distribution SR is the scene success rate on goals
    proposed by the frozen goal-proposing agent and accepted by the validity
    check, under the evaluation protocol of Section~\ref{sec:val_protocol};
    held-out SR is the 30-scene rate. Values are seed means with 95\,\%
    confidence intervals across the seeds listed. The single-agent models are
    omitted because their uniform training goal distribution coincides with the
    held-out suite, so their within-distribution and held-out rates are the same.}
    \label{tab:rq3}
\end{table}

The within-distribution evaluation separates the two effects that the headline
numbers confound. On goals the goal-proposing agent itself produces, the disc
model solves 61.4\,\% of scenes and the T-block model 11.7\,\%, across five seeds
each. The same models solve 7.3\,\% and 0.0\,\% of held-out scenes. The
single-agent models show no such separation: their uniform training goal
distribution coincides with the held-out suite, so their within-distribution and
held-out rates agree at roughly 79\,\%.

Two effects of comparable size are now measurable. The first is the coupled
gate. Within the proposer's own distribution, replacing the position-only disc
goal with the coupled T-block goal reduces scene success from 61.4\,\% to
11.7\,\%, a drop of roughly 50 percentage points. The second is the
generalisation collapse. The disc model loses a further 54 percentage points
between its within-distribution rate and the held-out suite, and the T-block
model falls from 11.7\,\% to zero. The coupled gate therefore bounds competence
on the proposer's own goals, while the narrow goal distribution prevents that
competence from transferring.

The per-axis training logs confirm that the coupled gate, not a missing skill,
limits the T-block. ASP-dPose satisfies the position gate on 41.0\,\% of its
training pushes and the rotation gate on 33.9\,\%, yet the combined gate is met
on only 11.7\,\% of within-distribution scenes. The two skills are present but
rarely co-satisfied. On the disc, where only position is demanded, the
within-distribution rate rises to 61.4\,\%, which confirms that the coupling is
what suppresses the T-block within-distribution. The residual transfer failure,
in which even the disc's competence does not reach the held-out suite, is a
property of the goal distribution itself. Alice proposes goals near her own final
pose, so the policy is trained on goals the held-out suite never presents, and
its competence is tied to those goals rather than to a transferable skill.

This answers RQ3. Self-play does not fail because a component of the method is
broken, and it does not fail from a single cause. Two effects compound: the
coupled multi-objective gate reduces the competence the policy acquires on its
own goals by roughly 50 percentage points, and the narrow goal distribution
prevents that competence from reaching the held-out suite. Both follow from the
same curriculum, and together they reduce the disc's 61.4\,\% to a held-out
7.3\,\% and the T-block's 11.7\,\% to zero.

\begin{figure}[H]
    \centering
    \includegraphics[width=\textwidth,height=0.42\textheight,keepaspectratio]{asp_diagnostic_2panel.png}
    \caption{The generalisation gap. Grey bars: within-distribution scene
    success rate (seed mean). Blue bars: held-out scene success rate with
    95\,\% confidence intervals. The single-agent models show no gap, because
    their training goal distribution coincides with the held-out suite; the
    self-play models collapse from their within-distribution rate to the held-out
    rate.}
    \label{fig:train_heldout}
\end{figure}

\subsection{Cross-Environment}
\label{sec:cross_env_results}

The second-simulator counterpart provides a single-seed qualitative check, and
each number below rests on that one seed. The numbers come from the primary
configuration of 64 environments and 15 pushes per rollout; the remaining
configurations of Section~\ref{sec:cross_env} are training-only. The
single-agent model solves 3.3\,\% of the 30 scenes, far below the 79.3\,\% of
the Isaac Lab model. The curriculum counterpart solves 9.0\,\%, and the
self-play counterpart solves 1.0\,\%. The lower absolute rates reflect the lower
fidelity of the second simulator and its single seed, not a contradiction of the
Isaac Lab result. What transfers is the ordering. The self-play model remains
the weakest of the three in both simulators. The curriculum model matches or
exceeds the plain single-agent model, consistent with the curriculum comparison
of Chapter~\ref{sec:conclusion}.

\subsection{Qualitative Behaviour}
\label{sec:qualitative}

Figure~\ref{fig:keyframes} shows keyframes of the single-agent PBRS policy
solving three held-out scenes. The policy approaches the block along the
object-relative offset, pushes along the decoded direction, and repeats until the
combined gate is met. Its failures, on the rotation-dominant scenes, show
position under control and orientation outside the 0.2\,rad tolerance, which
matches the per-type breakdown of Table~\ref{tab:main_results}.

\begin{figure}[H]
    \centering
    \includegraphics[width=0.32\textwidth]{rec_push_s11_key.png}\hfill
    \includegraphics[width=0.32\textwidth]{rec_push_s14_key.png}\hfill
    \includegraphics[width=0.32\textwidth]{rec_push_s21_key.png}
    \caption{Keyframes of the single-agent PBRS policy on three held-out scenes.}
    \label{fig:keyframes}
\end{figure}

The off-policy baseline, SAC with Hindsight Experience Replay, did not converge
within the available compute budget. Its status is recorded so that the on-policy
results are read against a concrete off-policy attempt rather than in isolation.
The failure is not evidence against the algorithm family: DDPG with hindsight
experience replay reaches a median success rate of 100\,\% on FetchPush, a
simulated tabletop pushing task, in \cite{Andrychowicz2017}. That published
result serves as plausibility calibration only. The simulator, the object, the
action space, and the compute budget all differ from the present task, so no
success rate is compared directly.
